\section{Introduction}\label{sec : intro}
Quantum groups, introduced by Drinfeld ~\cite{Dr85} and Jimbo~\cite{Ji85}, 
are $q$-deformations of universal enveloping algebras. 
In particular, when the Lie algebra is taken to be an affine Lie algebra, 
the corresponding quantum group is called a \emph{quantum affine algebra}. 
The study of finite-dimensional representations of quantum affine algebras 
has been an active area of research since the 1990s, 
owing to its significant applications not only in mathematics but also in physics, 
particularly in the theory of integrable systems and solvable lattice models
(see, for example, \cite{Baxter, FR92, DFJMN}). 
For a finite-dimensional simple Lie algebra $\mathfrak{g}$, 
let $\mathscr{C}_{\mathfrak{g}}$ be the category of finite-dimensional representations of 
the quantum affine algebra $U_q(\hat{\mathfrak{g}})$. 
In \cite{KKOP24}, Kashiwara-Kim-Oh-Park introduced a full subcategory 
$\mathscr{C}_{\mathfrak{g}}^{[a, b], \mathfrak{s}}$ 
of $\mathscr{C}_{\mathfrak{g}}$ and proved that it is a monoidal categorification of a cluster algebras, 
which is introduced in \cite{HL10}. 

Quiver Hecke algebras, introduced by Khovanov-Lauda and Rouquier~\cite{KL09, Ro08}, 
is infinite dimensional graded algebras associated with root data. 
Let $R$ be a symmetric quiver Hecke algebra over a field of characteristic $0$. 
Let $R$-gmod be a category of finite-dimensional graded module. 
It categorifies a quantum unipotent coordinate ring. 
For $w$ in the Weyl group $W$, 
the full subcategory $\mathscr{C}_{w}$ is introduced in \cite{KKKO18} and 
it is proved that $\mathscr{C}_{w}$ categorify a cluster algebra. 

Let $\mathscr{C}$ be $\mathscr{C}_{\mathfrak{g}}^{[a, b], \mathfrak{s}}$ or $\mathscr{C}_{w}$. 
Let $N = b-a+1$ for quantum affine case and $N = (\text{the length of }w)$ for quver Hecke case. 
In what follows, for $L, M \in \mathscr{C}$, even if $\mathscr{C} = \mathscr{C}_w$, 
we write the tensor product or the convolution product of  $L$ and $M$ as $L\otimes M$. 
Also, we ignore the grading shifts when $\mathscr{C} = \mathscr{C}_w$. 
For $L \in \Irr{\mathscr{C}}$, we call $L$ \emph{real} if $L\otimes L \in \Irr{\mathscr{C}}$ and 
\emph{imaginary} otherwise. 
The real and imaginary modules is really important issue to study the structure of 
the tensor product of two simple modules. 

To classify real and imaginary modules, the following conjecture is essential : 
\begin{Conj}
\label{Conj : real vs monomial intro}
    Let $\mathcal{C}$ be a monoidal categorification of a cluster algebra $\mathscr{A}$ 
    and let $\iota : K(\mathcal{C})_{\Z} \overset{\sim}{\to} \mathscr{A}$ be the corresponding 
    isomorphism. 
    Then, there is a bijection between the set of real simple modules 
    and the set of cluster monomials, i.e., 
    \[
    \{ \iota([L]) \mid L  \text{ is a real simple module in $\mathcal{C}$}\} 
    = \{ \text{cluster monomial in $\mathscr{A}$}\}. 
    \]
\end{Conj}
In this paper, we consider a monoidal analog of cluster-tilting theory and 
prove this conjecture when $\mathscr{C}$ is a monoidal categorification of a cluster algebra of acyclic type. 

Let $Q$ be a finite quiver and $\mathcal{C}_Q$ be the (generalized) cluster category. 
In the cluster-tilting theory, 
the following facts are known by \cite{BMRRT, IY08, DK08, FG} : 
\begin{enumerate}
  \item[(CT1)] For a basic cluster-tilting object $T$, $|T| = |Q_0|$, 
  \item[(CT2)] for an almost complete basic cluster-tilting object $\overline{T}$, 
  there exists exactly two indecomposable rigid objects which are complement of $\overline{T}$, 
  \item[(CT3)] $(M, M')$ is an exchange pair if and only if $\operatorname{Ext}^1_{\mathcal{C}_Q}(M, M') =1$, 
  \item[(CT4)] if $Q$ is acyclic, then any basic cluster-tilting object is reachable. 
\end{enumerate}
It is natural to ask that these fact also hold in the monoidal settings. 


First, we introduce \emph{maximal real set} in $\mathscr{C}$(Definition~\ref{Def : maximal real set}), 
which is a monoidal analog of a maximal rigid set. 
We call a family $\mathcal{M}$ of simple modules in $\mathscr{C}$ a maximal real set 
if it is a commuting family of real prime modules and it is maximal in such families. 
A maximal real set $\mathcal{M}$ is also a monoidal analog of cluster-tilting object 
by the following property(Proposition~\ref{Prop : simple modules commuting with a maximal real set}) : 
 if $X \in \Irr \mathscr{C}$ is commute with any element of $\mathcal{M}$, 
then $X$ is a monomial in $\mathcal{M}$. 
We prove that $\mathcal{M}$ is a maximal real set if and only if it is a Laurent family introduced in 
\cite{Laurentfamily} and prove the following corresponding to (CT1) : 
\begin{Thm}[Theorem~\ref{Thm : cardinality of maximal real sets}]
  Any maximal real set has exactly $N$ elements. 
\end{Thm}
By using properties of Laurent families in \cite{KK19, Laurentfamily}, 
we prove the following corresponding to (CT3) : 
\begin{Thm}[Theorem~\ref{Thm : characterization of exchange pairs}]
Let $L, L' \in \Irr\mathscr{C}$ be non-isomorphic real prime modules.
Then $(L, L')$ is an exchange pair if and only if
$\mathfrak{d}(L, L') = 1$.
\end{Thm}
For a maximal real set $\mathcal{M}$, 
we define the matrix $\widetilde B_{\mathcal M}$ using $\widetilde{\Lambda}$-invariant and 
prove the following corresponding to (CT2) : 
\begin{Thm}[Corollary~\ref{Cor : complete Lambda admissibility of maximal real sets}]
  A pair $(\mathcal{M}, \widetilde{B}_{\mathcal{M}})$ is a completely
$\Lambda$-admissible monoidal seed.
\end{Thm}
Also, for a real prime module $L \in \Irr\mathscr{C}$, 
we prove that $[L] \in K(\mathscr{C})_{\Z}$ is a prime element
(Proposition~\ref{Prop : real prime modules give prime elements}). 
By the property of $\mathscr{C}_{\mathfrak{g}}^{[a, b], \mathfrak{s}}$, 
we have the following : 
\begin{Thm}[Corollary~\ref{Cor : unique prime tensor factorization in Cg}]
For a real module $L \in \mathscr{C}_{\mathfrak{g}}$,
tensor factorization of $L$ into prime modules
is unique up to permutation.
\end{Thm}



For the cluster category $\mathcal{C}_{Q}$ and $\mathbf{g} \in \Z^Q$, 
there is a \emph{canonical decomposition} $\mathbf{g} = \mathbf{g}_1 \oplus \cdots\oplus \mathbf{g}_s$, 
which express the decomposition of generic presenintation in $\operatorname{PHom}(\mathbf{g})$ 
into indecomposable presentations. 
Viewing decomposition into indecomposable objects in a cluster category 
as the counterpart of prime factorization on the monoidal side, 
we want to transfer the structure of canonical decompositions to the monoidal setting. 
In the cluster category, $e$-invariant is introduced in \cite{DF15}. 
Since $\mathbf{g}=\mathbf{g}_1\oplus\cdots\oplus\mathbf{g}_s$ is the canonical decomposition 
if and only if $e(\mathbf{g}_i,\mathbf{g}_j)=0$ for all $i\neq j$, 
it is natural to compare the $\mathfrak d$-invariant and $e$-invariant. 
In \cite{Cao, Fei23}, 
the \emph{tropical $F$-invariant} is introduced and they compared the tropical $F$-invariant and 
$\mathfrak{d}, e$-invariants. 
In \cite{Cao}, it is proved that 
$\mathfrak{d}(L_1, L_2) = F_{\iota([L_1])}^t[(I\ 0)\deg^t\iota([L_2])]
+ F_{\iota([L_2])}^t[(I\ 0)\deg^t\iota([L_1])]$ if at least one of $L_1, L_2$ is reachable. 
We prove, without any assumptions, 
$\mathfrak{d}(L_1, L_2) \geq F_{\iota([L_1])}^t[(I\ 0)\deg^t\iota([L_2])]
+ F_{\iota([L_2])}^t[(I\ 0)\deg^t\iota([L_1])]$
 for $L_1, L_2 \in \Irr{\mathscr{C}}$
 (Theorem~\ref{Thm : tropical F invariants and Lambda inequalities}). 

Let consider the case that the cluster algebra which is categorified by $\mathscr{C}$ is acyclic type. 
For the case that  $Q$ is acyclic, 
it is proved in \cite{Fei23} that 
$F_{\mathbf{g}_1}^{\mathrm{gen}}[\mathbf{g}_2] = e(\mathbf{g}_1, \mathbf{g}_2)$ 
for $\mathbf{g}_1, \mathbf{g}_2\in \Z^{Q_0}$, 
where $F_{\mathbf{g}}^{\mathrm{gen}}$ is the generic $F$-polynomial. 
By comparing the Newton polytope of $F$-polynomial from simple module and generic $F$-polynomial, 
we obtain the following inequality : 
\begin{Thm}[Corollary~\ref{Cor : d dominates generic extensions}]
  Let $L_1, L_2 \in \Irr\mathscr{C}$.
Let $\widehat{\mathbf{g}}_i
\coloneq \deg^{t_1}\iota([L_i])$ and
$\mathbf{g}_i
\coloneq \operatorname{pr}_{\mathrm{uf}}\widehat{\mathbf{g}}_i$.
Then we have
\[
\mathfrak{d}(L_1, L_2)
\geq e(\mathbf{g}_1, \mathbf{g}_2)
+ e(\mathbf{g}_2, \mathbf{g}_1).
\]
\end{Thm}

As a result, if $\mathfrak{d} = 0$, then $e = 0$. 
By using this property, we prove the following corresponding to (CT4) : 
\begin{Thm}[Theorem~\ref{Thm : monoidal reachability for acyclic cluster algebras}]
Assume that the cluster algebra categorified by $\mathscr{C}$ is of acyclic type. 
For $L \in \Irr\mathscr{C}$,
$L$ is real if and only if $L$ is reachable.
Moreover, $L$ is real prime if and only if
$L$ corresponds to a cluster variable.
\end{Thm}




\subsection*{Acknowledgement}
I would like to express my sincere gratitude to my supervisor, 
Noriyuki Abe, for his continuous support and helpful discussions. 
Also I am deeply grateful to Hironori Oya for taking the time to discuss this work. 
This work was supported by JSPS KAKENHI Grant Number 26KJ0951. 

\subsection*{On the use of AI}
The author used Chat GPT to search related paper, check the draft and correctness of the argument, 
and improve the expression. 
All the ideas in this paper are the author's own. 
The author takes full responsibility for its content.





